
Une façon
équivalente de décrire ce dernier est la suivante : l'hypothèse que @@M0@@ est régulier implique que ces anneaux locaux sont factoriels (par le théorème de Auslander--Buchsbaum EGA IV 21.11.1), donc que tout faisceau inversible @@M1@@ sur @@M2@@ peut se prolonger en un faisceau inversible @@M3@@ sur @@M4@@ (EGA IV 21.6.11). Ceci posé, on vérifie aussitôt que l'image par (7.8.1) de @@M5@@ est @@M6@@, où on pose évidemment @@M7@@.

En fait, n'utilisant que l'hypothèse que les anneaux locaux de @@M8@@ sont factoriels et que @@M9@@ est localement intègre (à l'exclusion d'hypothèses de régularité et de séparation), on montre que @@M10@@ ne dépend pas du prolongement choisi @@M11@@ de @@M12@@ en un Module inversible sur @@M13@@, de sorte qu'on obtient ainsi directement un homomorphisme (7.8.1). Pour vérifier l'indépendance annoncée, on est ramené à prouver que si @@M14@@, alors @@M15@@. Or l'hypothèse signifie que @@M16@@, @@M17@@ un diviseur sur @@M18@@ de support @@M19@@ (cf. EGA IV 21.2.11), et on se réduit aussitôt au cas @@M20@@ \emph{intègre}, donc @@M21@@, mais alors @@M22@@ (@@M23@@ une uniformisante de @@M24@@), donc @@M25@@ et à fortiori @@M26@@.

\subsection*{7.9. Spécialisation et équivalence numérique.}
Supposons @@M27@@ propre, et @@M28@@ et @@M29@@ réguliers, et considérons les deux applications de spécialisation (7.4.1) et (7.7.4) :
@@M30@@
Du fait que l'homomorphisme
@@M31@@
